\chapter{Sustainable Development Approach}
This project aims to develop a product that brings us closer to a sustainable aviation future. In accordance with this, the team must perform its day-to-day operations sustainably, while also establishing a structured framework to minimize the environmental footprint of the final aircraft design.

\section{Operational Sustainability}

\begin{itemize}[leftmargin=*]
    \item \textbf{Digital-First Workflow and Resource Management:} The team will enforce a strict paperless policy for all internal communications, drafts, and data management. All deliverables, logbooks, and action items will be maintained digitally on the teams channel.
    
    \item \textbf{Low Carbon Footprint Commuting:} To minimize the project's direct carbon footprint, team members will rely on sustainable transport like bicycles, or shared public transport.

    \item \textbf{Minimize Computation:} Developing and validating design methods for aeroelastic phenomena requires significant computational power. To reduce energy consumption, the team will utilize efficient coding practices where possible. Code will be optimized and debugged locally before executing large-scale analytical workflows.
    
    \item \textbf{Sustainable Resource Allocation:} Office consumables and small materials provided by the Teaching Assistants will be utilized conservatively.
\end{itemize}

\section{Product Design and Lifecycle Sustainability}

Although the specific technical configuration will be finalized in subsequent project phases, the team will proactively embed sustainability into the core engineering decision-making process through the following framework:

\begin{itemize}[leftmargin=*]
    \item \textbf{Lifecycle Assessment (LCA) Integration:} As mandated by the project requirements, the team will conduct a comprehensive Lifecycle Assessment based on recent literature. This assessment will evaluate the environmental footprint of the UAV across its entire lifespan, from raw material extraction and manufacturing to operational use and eventual decommissioning.
    
    \item \textbf{Sustainability as a Trade-Off Criterion:} Environmental impact will serve as a weighted metric within formal trade-off analyses (e.g., using the Analytical Hierarchy Process). Material selection and manufacturing methods will be evaluated based on recyclability, toxicity, and embodied energy, carrying equal importance to standard aerospace metrics like yield strength and fatigue life.
    
    \item \textbf{SAF-Optimized Operations:} The VULCAN micro-turbine will operate exclusively on Sustainable Aviation Fuel (SAF). The team's aerodynamic and mission performance optimizations will prioritize overall aerodynamic efficiency and minimized fuel burn, ensuring the testbed itself operates with the lowest possible emission profile.
    
    \item \textbf{Modularity for Longevity:} The system is required to achieve a minimum operational life of 15 years. By designing a highly robust, modular airframe capable of supporting various interchangeable High-Aspect-Ratio (HAR) wings and tails, the team will future-proof the platform and drastically reduce the material waste associated with building entirely new test vehicles for future aeroelastic research.
\end{itemize}
